\begin{examplebox}
\textbf{例 1\quad 利用主圆秒求定距（最短应用）}

椭圆 $\frac{x^2}{4}+\frac{y^2}{3}=1$，任取切点 $P$；
左右焦点向切线作垂足 $H_1,H_2$，求 $OH_1,OH_2$。

\textbf{解}：$a=2$，由垂足共圆性质，
\[
\boxed{OH_1=OH_2=a=2}.
\]
\textbf{提示}：不是求 $PH_i$，而是求\textbf{椭圆中心到垂足}的距离，直接认出主圆。

\medskip
\textbf{例 2\quad 用长轴顶点一眼记住距离积}

椭圆 $\frac{x^2}{25}+\frac{y^2}{9}=1$，$a=5,b=3,c=4$。
在右顶点 $P=(5,0)$，切线为 $x=5$。
\[
d_1=5+4=9,\qquad d_2=5-4=1,
\qquad \boxed{d_1d_2=9=b^2}.
\]
同一个常数不因切点变化；但这里只是\textbf{记忆检查}，
不能用一个特例代替全椭圆证明。

\medskip
\textbf{例 3\quad 纠正“面积最大值恒为 $a^2$”}

椭圆 $\frac{x^2}{4}+\frac{y^2}{3}=1$，
$a=2,b=\sqrt3,c=1$。
由于 $c<b$，面积最大值应为
\[
\boxed{S_{\max}=2bc=2\sqrt3<4=a^2}.
\]
在上顶点 $P=(0,\sqrt3)$，切线 $y=\sqrt3$，
$H_1=(-1,\sqrt3)$、$H_2=(1,\sqrt3)$。
此时 $d_1=d_2=\sqrt3$，$H_1H_2=2$，
故 $S=(\sqrt3+\sqrt3)\cdot2/2=2\sqrt3$，确实取到最大值。
\textbf{迁移思路}：看到不等式上界，不仅要算上界，还要检查取等点是否存在。

\medskip
\textbf{例 4\quad 面积真能达到 $a^2$ 的完整构型}

仍取椭圆 $\frac{x^2}{25}+\frac{y^2}{9}=1$，
$c=4>b=3$。
考虑直线
\[
l:\ \sqrt7\,x+5y=20.
\]
由于椭圆沿法向量 $(\sqrt7,5)$ 的最大投影值为
$\sqrt{25\cdot7+9\cdot25}=20$，所以它是切线；
其切点为
\[
P=\left(\frac{5\sqrt7}{4},\frac94\right),\qquad
\frac{x_0^2}{25}+\frac{y_0^2}{9}=
\frac7{16}+\frac9{16}=1.
\]
两焦点到切线的距离为
\[
d_1=\frac{20+4\sqrt7}{\sqrt{32}},\qquad
 d_2=\frac{20-4\sqrt7}{\sqrt{32}}.
\]
因此 $d_1d_2=(400-112)/32=9$，且
\[
d_1+d_2=\frac{40}{\sqrt{32}}=5\sqrt2,
\qquad H_1H_2=\frac{8\cdot5}{\sqrt{32}}=5\sqrt2.
\]
两者恰好相等，所以
\[
\boxed{S=\frac12(5\sqrt2)(5\sqrt2)=25=a^2}.
\]
这不是只证明 ``$S\le25$''，而是\textbf{展示了真实的取等切线}。

\medskip
\textbf{例 5\quad 内外角平分线截距积（两条路线）}

椭圆 $\frac{x^2}{25}+\frac{y^2}{9}=1$，
取 $P=(3,12/5)$，则 $c=4$，且 $x_0y_0\ne0$。

\textbf{路线一（最快）}：
\[
x_M=\frac{16\cdot3}{25}=\frac{48}{25},\qquad
x_N=\frac{25}{3},\qquad
\boxed{x_Mx_N=16}.
\]

\textbf{路线二（从切线与法线独立算）}：
切线为
\[
\frac{3x}{25}+\frac{(12/5)y}{9}=1
\quad\Longleftrightarrow\quad 9x+20y=75.
\]
令 $y=0$，切线交点为 $N=(25/3,0)$。
法线方向 $(9,20)$，可写成
\[
(x,y)=\left(3,\frac{12}{5}\right)+t(9,20).
\]
令 $y=0$ 得 $t=-3/25$，故
$x_M=3-27/25=48/25$。
两条路线结果一致。

\medskip
\textbf{例 6\quad 光学性质与主圆的几何证明复现}

在椭圆上取非顶点 $P$，把 $F_1$ 关于切线翻折至 $Q$。
\begin{enumerate}
\item 反射性质保证 $F_2,P,Q$ 共线，且 $P$ 处于 $F_2,Q$ 之间；
因此 $F_2Q=PF_1+PF_2=2a$。
\item $H_1$ 为 $F_1Q$ 中点，$O$ 为 $F_1F_2$ 中点；
三角形 $F_1F_2Q$ 的中位线 $OH_1=\frac12F_2Q=a$。
\item 如果再取主圆上 $H_2$ 的对径点 $K=-H_2$，
则 $H_1,F_1,K$ 共线，点幂立得
$F_1H_1\cdot F_1K=a^2-c^2=b^2$。
\end{enumerate}
\textbf{迁移价值}：这三步可用在证明题，完全不用繁杂的坐标展开。

\medskip
\textbf{例 7\quad 特殊点检查}

对于 $\frac{x^2}{25}+\frac{y^2}{9}=1$：
\begin{itemize}
\item 上顶点 $P=(0,3)$：切线 $y=3$ 与 $x$ 轴平行；$M=O$，但 $N$ 不存在，所以不能写 $x_Mx_N=16$。
\item 右顶点 $P=(5,0)$：切线 $x=5$，法线是整条 $x$ 轴；焦点三角形退化，不能用通常的内外角平分线定理。距离积与垂足共圆仍然成立。
\end{itemize}
\end{examplebox}
